\section{Open Questions}
\label{sec:open-questions}

The following five questions are genuinely open after this replication --- each
is grounded in a specific gap in what we verified, and each has a concrete
next-step probe that could be run locally on free tooling.

\begin{enumerate}

\item \textbf{Modern arithmetic comparison.}
How does the Vedral--Barenco--Ekert ripple adder compare quantitatively to
modern quantum arithmetic primitives (Draper QFT adder, Cuccaro ripple adder,
Gidney's carry-lookahead) in gate count, T-count, and circuit depth at
$n{=}8, 16, 32$?
\emph{Basis:} the 1996 paper predates the Draper (2000) and Cuccaro (2004)
constructions; we reproduced VBE's $8n + \mathcal{O}(1)$ scaling with $R^{2}{=}1$
but never head-to-head'd it against those newer primitives.
\emph{Next steps:} implement Cuccaro and Draper under identical
NOT/CNOT/Toffoli accounting, run gate/depth counts on $n \in \{4,8,16,32\}$,
and report ratios.

\item \textbf{Cryptographic-scale memory formulas.}
Do the $7n{+}1 \to 4n{+}3$ memory formulas remain tight at $n \in \{1024, 2048\}$
(RSA-scale factoring) after Toffoli-to-Clifford+T compilation?
\emph{Basis:} we verified the formulas for $N \in \{15, 21, 35\}$
($n \le 6$) and reproduced the ``about 20 ions for $N{=}15$'' figure exactly,
but did not probe whether they hold after fault-tolerant compilation.
\emph{Next steps:} symbolically evaluate the formulas at $n \in \{512, 1024, 2048\}$,
then compile one mod-adder sub-block to Clifford+T (Toffoli $\to$ 7 T + 8 Clifford)
and report ancilla overhead and T-count.

\item \textbf{T-count competitiveness.}
Under a T-count / T-depth cost model (the relevant metric for fault-tolerant
quantum computation), is the VBE adder still competitive with Draper / Cuccaro,
or is the $8n$ Toffoli count a hidden cost that the paper's asymptotic
gate-count hides?
\emph{Basis:} the paper predates the modern insight that T-gates dominate
fault-tolerant arithmetic cost; our replication reports Toffoli+CNOT counts
(matching the paper) but did not decompose to $\{T, \text{Clifford}\}$.
\emph{Next steps:} decompose one full VBE adder at $n{=}8$ using the standard
7-T-per-Toffoli decomposition, count T-gates and T-depth, and compare against
Cuccaro's known $\sim\!14n$ T-count.

\item \textbf{ML-driven synthesis.}
Can an ML-driven synthesis tool (RL circuit optimization or a genetic-algorithm
ancilla-reduction pass) find a lower-gate-count reversible modular adder for
a fixed $N$ (say $N{=}15$ or $N{=}21$) than the VBE Fig.-4 construction?
\emph{Basis:} VBE's mod-adder (add-$a$, subtract-$N$, overflow-detect,
conditional add-back, temp reset) is provably correct but not proved optimal;
no modern ML synthesis pass has been run against these specific low-$n$
instances in this replication.
\emph{Next steps:} at $n{=}4$ ($N{=}15$), enumerate all functionally-equivalent
$\{\text{Toffoli}, \text{CNOT}\}$ circuits up to depth 30 via bounded SAT / RL
search; report whether any beats VBE while preserving the temp-reset invariant.

\item \textbf{Correctness extrapolation.}
Does our exhaustive basis-state validation ($n \le 4$ for plain adder,
$N \le 15$ for mod-adder) extrapolate to a proof of correctness at larger $n$,
or is there a known family of edge cases (all-ones inputs, near-$N$ inputs)
where implementation mistakes would only surface at $n \ge 8$?
\emph{Basis:} we verified only $n \in \{2,3,4\}$; VBE is inductive on $n$ so
$n{=}4$ correctness morally implies larger-$n$ correctness, but we ran no
SMT-based parameterized proof and did not target worst-case carry chains at
higher $n$.
\emph{Next steps:} run targeted large-$n$ worst-case sweeps at
$n \in \{8, 12\}$ ($a{=}b{=}N{-}1$, etc.), and encode the parameterized VBE
adder as a Z3 SMT formula to attempt a symbolic-$n$ correctness proof for the
CARRY/SUM primitives.

\end{enumerate}
